\documentclass[11pt]{article}
\usepackage[a4paper,margin=1in]{geometry}
\usepackage{amsmath,amssymb,bm}
\usepackage{listings}
\usepackage{xcolor}
\usepackage{hyperref}

\hypersetup{colorlinks=true,linkcolor=blue!50!black,urlcolor=blue!50!black}

\lstset{
  basicstyle=\ttfamily\small,
  breaklines=true,
  columns=fullflexible,
  frame=single,
  framerule=0.3pt,
  rulecolor=\color{gray!50},
  xleftmargin=1em,xrightmargin=1em,
  aboveskip=4pt,belowskip=4pt
}

\newcommand{\F}{F_{\!a}}            % Faraday
\newcommand{\Rg}{R_{\!g}}           % gas constant
\newcommand{\eps}{\varepsilon}
\newcommand{\dV}{\,\mathrm{d}V}
\newcommand{\dS}{\,\mathrm{d}S}
\newcommand{\nvec}{\mathbf{n}}
\newcommand{\Eu}{\mathbf{E}}
\newcommand{\sgn}{\mathrm{sgn}}

% Layout helpers: each entry = three labelled blocks (Inputs / PDE / Code / Notes)
\newenvironment{ecmEntry}[1]{%
  \subsection{#1}%
}{}
\newcommand{\inputsRow}[1]{\par\noindent\textbf{Inputs.} #1\par}
\newcommand{\pdeRow}[1]{\par\noindent\textbf{PDE / weak form.}\par #1\par}
\newcommand{\codeRow}{\par\noindent\textbf{Residual (verbatim).}\par}
\newcommand{\notesRow}[1]{\par\noindent\textbf{Notes.} #1\par}

\title{ECM\_test: Physics and PDE Catalog\\
       \large Per-class documentation of every MOOSE object in \texttt{ecm\_test/src/}}
\author{}
\date{\today}

\begin{document}
\maketitle
\tableofcontents

\section*{Conventions}
\addcontentsline{toc}{section}{Conventions}
\begin{description}
  \item[Variables.] $u$ is the unknown solved by a given residual; $c,\,\phi,\,T$ denote
    concentration, electric/chemical potential, and temperature when those couplings are explicit.
  \item[Constants.] $\F\!=\!96485\,\mathrm{C\,mol^{-1}}$ Faraday, $\Rg\!=\!8.3145\,\mathrm{J\,mol^{-1}K^{-1}}$ gas constant,
    $T$ temperature, $z$ charge number. Several classes use the shorthand
    $F_{RT}\!\equiv\!F/(RT)$ (so $F_{RT}\!\approx\!38.68\,\mathrm{V^{-1}}$ at 300\,K).
  \item[Discrete weak form.] Each MOOSE residual is the integral
    $\int_\Omega r(u,\nabla u)\,\psi\,\dV$ for kernels, or $\int_{\partial\Omega} r\,\psi\dS$
    for IntegratedBCs/InterfaceKernels, where $\psi$ is the test function. We write the
    equation in strong form and quote the precise C++ residual.
  \item[Hard-coded scaling.] Several classes apply unit-conversion prefactors of $10$, $100$,
    $10\,000$, etc. These are reproduced verbatim and flagged in \emph{Notes}.
  \item[Coupling kernels.] When a class advertises a Butler--Volmer (BV) flux, the
    overpotential argument depends on which boundary it sits on; we write the activation
    argument as $\Delta\Phi$ and define it in each entry.
\end{description}

\section*{Reading guide}
\addcontentsline{toc}{section}{Reading guide}
For each class we give: the inputs (coupled variables, material properties, scalar parameters),
the strong-form PDE term it discretises, the verbatim C++ residual returned by
\texttt{computeQpResidual}/\texttt{computeQpProperties}/etc., and brief notes on quirks,
non-obvious choices, or potential issues. Quirks include duplicated implementations,
hard-coded magic factors, missing test-function multipliers, latent bugs we spotted while
reading, and AD-vs-non-AD legacy duplicates.

%==========================================================================
\section{Bulk Kernels (\texttt{src/kernels/})}
%==========================================================================

\subsection*{1. Butler--Volmer source family}

\begin{ecmEntry}{ADButlerVolmerForce}
\inputsRow{material property \texttt{mat\_prop\_name} (typically
\texttt{butler\_volmer\_current\_force} or \texttt{\dots\_flux\_force} from
\hyperref[bv-mat]{\textsc{ADButlerVolmerMaterial}}); \texttt{scale}; \texttt{base\_name}.}
\pdeRow{Volumetric reaction source
$r(\mathbf{x}) = \texttt{scale}\cdot p_{\text{mat}}(\mathbf{x})$, applied
to whichever equation $u$ this kernel decorates (typically the Li or electron equation):
\[
  \int_\Omega \texttt{scale}\,p_{\text{mat}}\,\psi \dV.
\]}
\codeRow\begin{lstlisting}
return _scale * _mat_property[_qp] * _test[_i][_qp];
\end{lstlisting}
\notesRow{Lifts a precomputed BV scalar onto the residual; the actual BV physics lives
in \textsc{ADButlerVolmerMaterial}.}
\end{ecmEntry}

\begin{ecmEntry}{ADCoupledButlerVolmerForce}
\inputsRow{coupled $\phi_{\text{electrolyte}},\,\phi_{\text{electrode}}$;
\texttt{exchange\_current\_density} $i_0$; equilibrium potential $U_{eq}$ (material);
$\F,T,\Rg$.}
\pdeRow{Symmetric ($\alpha\!=\!\tfrac{1}{2}$) Butler--Volmer reaction term:
\[
  r = 2\,i_0\,\sinh\!\bigl(\tfrac{1}{2}\eta/(\Rg T/\F)\bigr),
  \qquad \eta = \phi_{\text{electrode}} - \phi_{\text{electrolyte}} - U_{eq}.
\]}
\codeRow\begin{lstlisting}
RTF = R*T/F;
eta = phi_e - phi_l - Ueq;
return 2*i0*std::sinh(0.5*eta/RTF);
\end{lstlisting}
\notesRow{The residual omits the explicit $\psi$ multiplier present in
\textsc{ADButlerVolmerForce}; in MOOSE this means the value is treated as a strong-form
source. Verify against intended kernel base if accuracy is critical.}
\end{ecmEntry}

\begin{ecmEntry}{ADMultipleButlerVolmerForce}
\inputsRow{vector \texttt{materials} of \textsc{ADButlerVolmerMaterial} instances;
\texttt{mat\_prop\_name}; \texttt{scale}.}
\pdeRow{Sums BV contributions from a list of materials (e.g.\ multi-particle electrodes):
\[
  r = \texttt{scale}\sum_{m=1}^{M} p^{(m)}_{\text{mat}}(\mathbf{x}),\qquad
  \int_\Omega r\,\psi\,\dV.
\]}
\codeRow\begin{lstlisting}
res = 0;
for (auto m : _bv_materials)
  res += _scale * m->_mat_prop[_qp];
return res * _test[_i][_qp];
\end{lstlisting}
\notesRow{Use when one finite element overlaps multiple active material phases;
prefactors are summed before integration.}
\end{ecmEntry}

\subsection*{2. Diffusion of charged / chemo-mechanically coupled species}

\begin{ecmEntry}{ADChemoMechanoDiffusionBase (templated)}
\label{kernel-cmm-diff}
\inputsRow{coupled \texttt{stress\_based\_chemical\_potential} $\mu_\sigma$;
$\Rg=8.314426,\,T=298$; diffusivity $D$ (scalar or tensor) inherited from
\texttt{ADMatDiffusionBase}.}
\pdeRow{Adds a chemo-mechanical drift to the Fickian flux:
\[
  -\nabla\cdot\mathbf{J} = -\nabla\cdot\!\Bigl(D\,\nabla u + \frac{D u}{\Rg T}\nabla\mu_\sigma\Bigr).
\]
The residual returned is $\mathbf{J}_{\text{precomp}}$ to be dotted with $\nabla\psi$.}
\codeRow\begin{lstlisting}
auto residual = ADMatDiffusionBase<T>::precomputeQpResidual();
if (_mu_coupled)
    residual += _diffusivity[_qp] * _u[_qp] * (*_grad_mu)[_qp]
              / (_gas_constant * _temperature);
return residual;
\end{lstlisting}
\notesRow{Sets \texttt{use\_displaced\_mesh = false} (kinematics handled by the swelling
material). Two thin wrappers \textsc{ADChemoMechanoDiffusion} (scalar $D$) and
\textsc{ADChemoMechanoAnsioDiffusion} (tensor $D$) instantiate the templated base.}
\end{ecmEntry}

\begin{ecmEntry}{ADStressBasedChemicalPotential}
\inputsRow{$\mu_{\text{swell}},\,\mu_{\text{elastic}}$ (AD material properties produced by
the swelling base \S\ref{mat-swelling-base}).}
\pdeRow{Algebraic constraint defining a stress-based chemical-potential auxiliary variable:
\[
  u - \mu_{\text{swell}} - \mu_{\text{elastic}} = 0.
\]}
\codeRow\begin{lstlisting}
return (_u[_qp] - _swelling_chem_pot[_qp] - _elastic_chem_pot[_qp]) * _test[_i][_qp];
\end{lstlisting}
\notesRow{Pairs with \hyperref[kernel-cmm-diff]{\textsc{ADChemoMechanoDiffusion}}: the
diffusion kernel reads $\nabla u$ as $\nabla\mu_\sigma$. \texttt{use\_displaced\_mesh=false}.}
\end{ecmEntry}

\begin{ecmEntry}{ADChargedTransportAnisotropy / ChargedTransport}
\inputsRow{tensor (resp.\ scalar) diffusivity $D$ as AD material; hard-coded factor $10$.}
\pdeRow{Pure Laplacian with anisotropic conductivity and a unit-conversion factor:
\[
  -\nabla\cdot(10\,\bm{D}\,\nabla u),\qquad
  -\nabla\cdot(10\,D\nabla u).
\]}
\codeRow\begin{lstlisting}
return 10.0 * raw_value(_diffusivity_coef[_qp]) * _grad_test[_i][_qp] * _grad_u[_qp];
\end{lstlisting}
\notesRow{Both classes call \texttt{raw\_value} on the diffusivity, which strips AD
sensitivities; the Jacobian therefore ignores derivatives through $D$.}
\end{ecmEntry}

\begin{ecmEntry}{ADMigrationDiffusion}
\inputsRow{coupled $c$; conductivity $\sigma$ (material); $c_0,\,z=1,\,F_{RT}=36.41,\,$\texttt{scale}.}
\pdeRow{Combined diffusion + migration current using log-concentration linearisation,
applied as a current-equation source:
\[
  -\nabla\cdot\!\Bigl(\frac{\sigma}{F_{RT}}\bigl(\tfrac{c_0}{c}\nabla c + z\,F_{RT}\nabla u\bigr)\Bigr).
\]}
\codeRow\begin{lstlisting}
term1 = c0/c * _grad_c;
term2 = F_RT * _grad_u;
return scale * sigma/F_RT * _grad_test * (term1 + zIons*term2);
\end{lstlisting}
\notesRow{AD analogue of the legacy \textsc{Migration} kernel; conventions and constants
differ from \textsc{NernstPlanckConvection}.}
\end{ecmEntry}

\begin{ecmEntry}{ADNernstPlanckConvection}
\inputsRow{coupled potential $V$; diffusivity $D$ (material); $z=1,\,F_{RT}=38.68,\,$\texttt{scale}.}
\pdeRow{Fick + drift Nernst--Planck flux, in mixed convective/diffusive form:
\[
  -\nabla\cdot(D\nabla u) - z F_{RT}\,D\,\nabla V\!\cdot\!\nabla u.
\]}
\codeRow\begin{lstlisting}
return D*_grad_test*_grad_u + D*zIons*F_RT*_grad_V*_test*_grad_u;
\end{lstlisting}
\notesRow{Migration term uses $\psi\,\nabla V\!\cdot\!\nabla u$ (gradient-of-$u$ contracted
with $\nabla V$) instead of $\nabla\psi\cdot c\nabla V$; this is a convective rather
than divergence form.}
\end{ecmEntry}

\begin{ecmEntry}{NernstPlanckConvection (legacy)}
\inputsRow{coupled $V$; $D$ (material); $z=1,\,F_{RT}=38.68$.}
\pdeRow{$\;-\,z\,(F_{RT} D\,\nabla V)\cdot\nabla u\;$ (convective form, no diffusion piece).}
\codeRow\begin{lstlisting}
ion_velocity = F_RT * D * _grad_V;
return -_test * zIons * ion_velocity * _grad_u;
\end{lstlisting}
\notesRow{Provides hand-coded Jacobians wrt $u$ and $V$. Uses
\texttt{Kernel} (non-AD) base.}
\end{ecmEntry}

\begin{ecmEntry}{NernstPlanckConvectionScaled}
\inputsRow{coupled $V$; $D$; $z=1,\,F_{RT}=38.68,\,$\texttt{scale}.}
\pdeRow{Scaled twin of the previous entry: $\;-\,$\texttt{scale}$\,z\,F_{RT}D\,\nabla V\cdot\nabla u$.}
\codeRow\begin{lstlisting}
ion_velocity = F_RT * D * _grad_V;
return -_test * scale * zIons * ion_velocity * _grad_u;
\end{lstlisting}
\notesRow{Identical residual modulo the explicit \texttt{scale}.}
\end{ecmEntry}

\begin{ecmEntry}{NernstPlanckConvectionEtime}
\inputsRow{coupled scalar electric field magnitude \texttt{Efield}; $D$; $z=1,\,F_{RT}=38.68,\,$\texttt{scale}.}
\pdeRow{Drift-only term using the field directly (not $\nabla V$):
$\;-\,$\texttt{scale}$\,z\,F_{RT}\,D\,u\,E$.}
\codeRow\begin{lstlisting}
return -scale * _test * zIons * F_RT * D * _u * _Efield;
\end{lstlisting}
\notesRow{For time-domain electromagnetic coupling where $E$ is a primary variable.}
\end{ecmEntry}

\subsection*{3. Poisson and current-flux kernels}

\begin{ecmEntry}{ADPoissonEquation}
\inputsRow{coupled concentration $c$; $\F=96485.34$, $\sigma$, $z$, \texttt{scale}.}
\pdeRow{$\;-\nabla\cdot\!\bigl((\sigma/\F)\,\nabla u\bigr) = z\,c\;$ (Poisson with charge source).}
\codeRow\begin{lstlisting}
return (sigma/F) * scale * _grad_test * _grad_u
       - _test * z * c;
\end{lstlisting}
\notesRow{Filename has typo (\texttt{ADPoissionEquation.C}); class name is correct.
$\sigma$ and $z$ are scalar \texttt{Real}, not material properties.}
\end{ecmEntry}

\begin{ecmEntry}{ChargeDensity}
\inputsRow{coupled $c_{\text{ions}}$; $z=1,\,$\texttt{scale}=1.}
\pdeRow{Source for the Poisson equation: $\;r=-\,$\texttt{scale}$\,z\,c$.}
\codeRow\begin{lstlisting}
return -scale * _test * zIons * _conIons;
\end{lstlisting}
\notesRow{Pairs naturally with \textsc{ADPoissonEquation}.}
\end{ecmEntry}

\begin{ecmEntry}{ADCurrentFluxEEL}
\inputsRow{coupled $c$; $\sigma$ (material); $\F=964523,\,\Rg=8.845,\,T=298,\,c_0,\,$\texttt{scale}.}
\pdeRow{Dilute-solution current density from the EEL formulation:
\[
  \mathbf{i} = \sigma\Bigl(\nabla u + \frac{1}{\F}\nabla\mu\Bigr),\quad
  \nabla\mu = \frac{\Rg T c_0}{c}\nabla c.
\]}
\codeRow\begin{lstlisting}
gradMu  = R*T*c0/c * _grad_c;
gradPhi = _grad_u;
return _scale * sigma * _grad_test * (gradPhi + gradMu/F);
\end{lstlisting}
\notesRow{Default \texttt{faraday=964523} differs from the standard $96485$; verify before
trusting. Imported from the \texttt{eel} parent app.}
\end{ecmEntry}

\begin{ecmEntry}{ADMassFluxEEL}
\inputsRow{coupled $V$; mobility $b$, conductivity $\sigma$ (materials);
$\F=96485,\,\Rg=8.854,\,T=298,\,c_0,\,$\texttt{scale}.}
\pdeRow{Nernst--Planck mass flux for Li$^+$:
$\;\mathbf{J}_c = b\,\Rg T c_0\,\nabla u - (\sigma/\F)\nabla V.$}
\codeRow\begin{lstlisting}
flux_c   = mobility * R*T*c0 * _grad_u / _test;
flux_phi = -conductivity/F * _grad_V;
return _scale * (flux_c + flux_phi) * _grad_test;
\end{lstlisting}
\notesRow{The division by \texttt{\_test} followed by multiplication by \texttt{\_grad\_test}
is unusual; should be reviewed against the original eel formulation.}
\end{ecmEntry}

\begin{ecmEntry}{ADThermalFlux}
\inputsRow{conductivity $k$ (material), \texttt{scale}.}
\pdeRow{Fourier conduction $-\nabla\cdot(k\,\nabla u)$.}
\codeRow\begin{lstlisting}
return k * scale * _grad_test * _grad_u;
\end{lstlisting}
\notesRow{Uses non-AD \texttt{getMaterialProperty} so AD sensitivities through $k$ are lost.}
\end{ecmEntry}

\begin{ecmEntry}{ParamDiffusion}
\inputsRow{conductivity $\sigma$ (AD material).}
\pdeRow{$\;-\nabla\cdot(10\,000\,\sigma\,\nabla u)\;$ (Laplacian with hard-coded
unit-conversion factor).}
\codeRow\begin{lstlisting}
return 10000 * conductivity * _grad_test * _grad_u;
\end{lstlisting}
\notesRow{Likely a development precursor to \textsc{ChargedTransport}; magic factor
$10\,000$ is suspicious.}
\end{ecmEntry}

\begin{ecmEntry}{Test}
\inputsRow{$A=1$.}
\pdeRow{$\;-\nabla\cdot(A\,\nabla u)\;$ (sanity-check Laplacian).}
\codeRow\begin{lstlisting}
return A * _grad_test * _grad_u;
\end{lstlisting}
\notesRow{Development stub. Not used in any production input.}
\end{ecmEntry}

\subsection*{4. Legacy non-AD electrochemistry trio}

These three kernels share a \texttt{DerivativeMaterialInterface} contract: they read
parsed-derivative materials and assemble manual Jacobians.

\begin{ecmEntry}{Conduction}
\inputsRow{coupled $c_p,c_v$; \texttt{Q\_name}, \texttt{QM\_name} as materials with
parsed derivatives wrt $c_p,c_v$.}
\pdeRow{$\;\nabla\psi\cdot(Q\,\nabla u + Q_M\,\nabla c_p)\;$.}
\codeRow\begin{lstlisting}
return Q * _grad_u * _grad_test + QM * _grad_cp * _grad_test;
\end{lstlisting}
\notesRow{Hand-coded \texttt{computeQpJacobian}/\texttt{computeQpOffDiagJacobian} use the
parsed derivatives.}
\end{ecmEntry}

\begin{ecmEntry}{Migration (legacy)}
\inputsRow{coupled $c_p,c_v$; \texttt{Q\_name}, \texttt{QM\_name}.}
\pdeRow{$\;\nabla\psi\cdot(Q\,\nabla c_v + Q_M\,\nabla c_p)\;$.}
\codeRow\begin{lstlisting}
return Q * _grad_cv * _grad_test + QM * _grad_cp * _grad_test;
\end{lstlisting}
\notesRow{Identical structure to \textsc{Conduction} but uses $\nabla c_v$ instead of
$\nabla u$; intended for the dual variable in coupled chemo-electric models.}
\end{ecmEntry}

\begin{ecmEntry}{Kinetics}
\inputsRow{coupled $c_p,c_v$; parsed material $F$.}
\pdeRow{Source $r=F$ (parsed free-energy or reaction-rate expression).}
\codeRow\begin{lstlisting}
return _F * _test;
\end{lstlisting}
\notesRow{Off-diagonal Jacobian uses derivatives of $F$ wrt $u,\,c_p,\,c_v$.}
\end{ecmEntry}

\begin{ecmEntry}{SourceTerm (legacy)}
\inputsRow{material $a$ (\texttt{srcCoef}); user \texttt{Function} $f$; $F_{RT}=38.68$.}
\pdeRow{Nonlinear quadratic source: $\;r=F_{RT}\,a\,u^2 - f(t,\mathbf{x})$.}
\codeRow\begin{lstlisting}
return _test * (F_RT * srcCoef * u*u - f(t,x));
\end{lstlisting}
\notesRow{Used in \texttt{CHEMICAL\_REACTIONS}-style problems.}
\end{ecmEntry}

\subsection*{5. Phase-field family}

\begin{ecmEntry}{FreeEnergyDouble}
\inputsRow{$A$ (AD material); \texttt{scale}=1.}
\pdeRow{Derivative of the standard double-well bulk energy
$f(u)=A u^2(1-u)^2$:
\[
  r = \frac{\partial f}{\partial u} = 2A(1-u)^2 u - 2A(1-u)u^2 = 2Au(1-u)(1-2u).
\]}
\codeRow\begin{lstlisting}
return scale * _test * (2*A*pow(1-u,2)*u - 2*A*(1-u)*pow(u,2));
\end{lstlisting}
\notesRow{Standard Cahn--Hilliard / Allen--Cahn bulk term.}
\end{ecmEntry}

\begin{ecmEntry}{ChemicalEnergy}
\inputsRow{$w$ (AD material); \texttt{scale}=1.}
\pdeRow{Derivative of $f(u)=w\,(u-1)^2 u^2$ scaled by $30$:
$r=30\,w(u-1)^2 u^2 \,\psi$ \emph{(per the source — note this is the integrand,
not the derivative; see Notes)}.}
\codeRow\begin{lstlisting}
return _test * scale * w * 30 * (u-1)*(u-1) * u*u;
\end{lstlisting}
\notesRow{The form returned is $f$ itself, not $\partial f/\partial u$. If used as a
phase-field bulk-energy gradient kernel this is mathematically inconsistent; verify
intent against input files.}
\end{ecmEntry}

\begin{ecmEntry}{NucleationEnergy}
\inputsRow{$\Delta G,\,\gamma$ (AD materials); \texttt{scale}=1.}
\pdeRow{Linear penalty driving $u\to 1$:
$\;r = 2\,$\texttt{scale}$\,\Delta G\,\gamma\,(u-1)$.}
\codeRow\begin{lstlisting}
return 2 * _test * scale * dG * gamma * (u - 1);
\end{lstlisting}
\notesRow{Models the energy barrier to nucleation; positive $\Delta G\,\gamma$ stabilises
$u=1$ (nucleated phase).}
\end{ecmEntry}

\begin{ecmEntry}{CouplingFreeEnergy}
\inputsRow{coupled order parameter $v$; $c$ (AD material); \texttt{scale}=1.}
\pdeRow{Derivative wrt $u$ of the cross-energy $f(u,v)=c\,u^2 v^2$:
$\;\partial f/\partial u = 2 c\,u\,v^2$.}
\codeRow\begin{lstlisting}
return 2 * scale * _test * c * _u * coupledVar*coupledVar;
\end{lstlisting}
\notesRow{Used in multi-phase free-energy decompositions.}
\end{ecmEntry}

\begin{ecmEntry}{CoupledTimeDerivativePhase}
\inputsRow{coupled order parameter $\phi$; $n=1,\,F=1,\,$\texttt{scale}=1; $C$ (AD material).}
\pdeRow{Time derivative $r = n\F\,$\texttt{scale}$\,C\,\dot{\phi}$ (couples a Poisson-type
equation to phase-field reaction kinetics).}
\codeRow\begin{lstlisting}
return n * F * scale * C * _phi_var_dot * _test;
\end{lstlisting}
\notesRow{Often paired with \textsc{ElectrodeDrivingForce} below.}
\end{ecmEntry}

\begin{ecmEntry}{ElectrodeDrivingForce}
\inputsRow{coupled $\phi,\,\phi_{\text{ref}},\,c$; $h$ (AD material, interpolation function);
$\alpha=\beta=1,\,n=1,\,\Rg=8.84,\,T=298,\,\F=96485,\,$\texttt{scale}=1.}
\pdeRow{Phase-field BV-style driving force:
\[
  r = \texttt{scale}\,h\bigl[\exp(\alpha\,\Delta G) - \exp(-\beta\,\Delta G)\bigr],\quad
  \Delta G = \frac{n\F(\phi-\phi_{\text{ref}})}{\Rg T}.
\]}
\codeRow\begin{lstlisting}
dG = n*F*(pot - ref_pot) / (R*T);
return _test * scale * h * (exp(alpha*dG) - exp(-beta*dG));
\end{lstlisting}
\notesRow{No exchange-current $i_0$ prefactor (set externally via $h$). The coupled
\texttt{conc} parameter is declared but unused in the residual.}
\end{ecmEntry}

\begin{ecmEntry}{PhaseFieldLaplace}
\inputsRow{$k_0,\,$\texttt{scale} (AD materials).}
\pdeRow{Laplacian with two material prefactors: $\;-\nabla\cdot(k_0\,$\texttt{scale}$\,\nabla u)$.}
\codeRow\begin{lstlisting}
return k0 * scale * _grad_test * _grad_u;
\end{lstlisting}
\notesRow{Both prefactors fetched as \texttt{ADMaterialProperty<Real>}.}
\end{ecmEntry}

\subsection*{6. Sub-particle approximation}

\begin{ecmEntry}{SurfaceConcentration2ParameterAppox}
\inputsRow{material flux $j$, $D$ (AD); coupled average concentration $\bar c$;
particle radius $R$; \texttt{scale}.}
\pdeRow{Bernardi--Newman 2-term polynomial approximation (\texttt{ADKernelValue},
adds the value not its divergence):
\[
  r = \bar c - u - \frac{R}{5\,\mathrm{scale}\,D}\,j.
\]
The kernel pins $u = c_{\text{surf}} \approx \bar c + R\,j/(5D)$.}
\codeRow\begin{lstlisting}
return -flux/5/D * particle_size * scale + avg_c - u;
\end{lstlisting}
\notesRow{Used to avoid resolving radial sub-particle diffusion; valid in the
diffusion-limited regime.}
\end{ecmEntry}

%==========================================================================
\section{Boundary Conditions (\texttt{src/bcs/})}
%==========================================================================

The BV BCs share a common skeleton: an exchange current $i_{\text{ex}}$ (an AD material
property) modulated by the symmetric/asymmetric activation
$k_1 = \exp(\alpha\,F_{RT}\,\Delta\Phi)$, $k_2=\exp(-(1-\alpha)\,F_{RT}\,\Delta\Phi)$, with
$\Delta\Phi$ depending on which interface and which equation. We give $\Delta\Phi$ explicitly
for each entry. Several BCs add an ohmic flux $-\sigma\nabla\phi\cdot\nvec$ alongside.

\subsection*{1. Generic flux and Dirichlet}

\begin{ecmEntry}{MaterialDiffusionFluxBC}
\inputsRow{$D$ (AD material).}
\pdeRow{$\;-D\,\nabla u\cdot\nvec\;$ on $\partial\Omega$.}
\codeRow\begin{lstlisting}
return -D * _grad_u * _normals[_qp] * _test[_i][_qp];
\end{lstlisting}
\notesRow{Standard outward-normal Fickian-flux BC.}
\end{ecmEntry}

\begin{ecmEntry}{ScaledCoupledVarNeumannBC}
\inputsRow{coupled $v$; \texttt{scale}.}
\pdeRow{Imposes $\;\partial u/\partial n = $\,\texttt{scale}$\,v\;$ via residual
$-$\,\texttt{scale}$\,v\,\psi$.}
\codeRow\begin{lstlisting}
return -_test[_i][_qp] * v * scale;
\end{lstlisting}
\end{ecmEntry}

\begin{ecmEntry}{PorousFlowSinkScaledCoupledVar}
\inputsRow{coupled $v$; \texttt{scale}; inherits MOOSE \texttt{PorousFlowSink}.}
\pdeRow{Multiplies an existing porous-flow sink residual by $v\,$\texttt{scale}.}
\codeRow\begin{lstlisting}
return PorousFlowSink::computeQpResidual() * v * scale;
\end{lstlisting}
\end{ecmEntry}

\begin{ecmEntry}{DirichletEn / DirichletLi}
\inputsRow{coupled scalar (\texttt{potLi} or \texttt{potEn}, indexed at element 0).}
\pdeRow{Nodal BC pinning $u$ to a global scalar:\\
\textsc{DirichletEn}: $u - \texttt{potLi}[0]=0$;\quad
\textsc{DirichletLi}: $u - \texttt{potEn}[0]=0$.}
\codeRow\begin{lstlisting}
return _u[_qp] - coupled_scalar[0];
\end{lstlisting}
\notesRow{Treats the coupled variable as a single global scalar (e.g.\ a battery cell
voltage).}
\end{ecmEntry}

\begin{ecmEntry}{SingleSEElectrodeNeumann}
\inputsRow{\texttt{applied\_current} (AD material).}
\pdeRow{$\;-i_{\text{applied}}\;$ on $\partial\Omega$.}
\codeRow\begin{lstlisting}
return -_test[_i][_qp] * _applied_current[_qp];
\end{lstlisting}
\end{ecmEntry}

\subsection*{2. Butler--Volmer boundary kinetics}

\begin{ecmEntry}{ADButlerVolmerBC}
\inputsRow{$i_0,\,\F=96485.33,\,\Rg=8.3145,\,T=298$; constant or function applied current
$i$; $U_{eq}$ (AD material).}
\pdeRow{Symmetric BV plus applied current at the boundary:
\[
  r = (i - i_0\,(e^{\Delta\Phi} - e^{-\Delta\Phi}))\psi,\quad
  \Delta\Phi = \tfrac{1}{2}\F(u-U_{eq})/(\Rg T).
\]}
\codeRow\begin{lstlisting}
prefac = 0.5*F/(R*T);
dPhi   = prefac*u; if (Ueq) dPhi -= prefac*Ueq;
fac    = exp(dPhi) - exp(-dPhi);
i      = (current_param ? current_param : current_function);
return (i - i0 * fac) * _test;
\end{lstlisting}
\end{ecmEntry}

\begin{ecmEntry}{ButlerVolmerIonics}
\inputsRow{$F_{RT}=0.03868\;\mathrm{mV^{-1}}$; $\alpha=0.5$; $V_{\text{Li,ref}}$;
$i_{\text{ex}}$.}
\pdeRow{BV across an ionic SE \,\textbar\, Li-metal interface:
\[
  r = -i_{\text{ex}}(k_1-k_2)\,\psi,\quad
  \Delta\Phi = V_{\text{Li,ref}}-u.
\]}
\codeRow\begin{lstlisting}
k1 = exp( alpha*F_RT*(LiPotRef - u) );
k2 = exp(-(1-alpha)*F_RT*(LiPotRef - u) );
return -_test * i_ex * (k1 - k2);
\end{lstlisting}
\notesRow{$F_{RT}$ in mV$^{-1}$ — confirm potentials are in mV when wired up.}
\end{ecmEntry}

\begin{ecmEntry}{ButlerVolmerMiec}
\inputsRow{coupled $\phi_e$; $F_{RT},\,\alpha$, $V_{\text{Li,ref}},\,i_{\text{Li,ref}},\,i_{\text{ex}}$.}
\pdeRow{MIEC variant including the electron potential in the activation argument; rectifies
to disable cathodic branch:
\[
  \Delta\Phi = V_{\text{Li,ref}}-u-\phi_e,\quad
  r = \begin{cases} -\bigl(i_{\text{Li,ref}}-i_{\text{ex}}(k_1-k_2)\bigr)\psi, & k_1<k_2,\\
                     -i_{\text{Li,ref}}\psi, & k_1\ge k_2.\end{cases}
\]}
\codeRow\begin{lstlisting}
k1 = exp( alpha*F_RT*(LiPotRef - u - potEn) );
k2 = exp(-(1-alpha)*F_RT*(LiPotRef - u - potEn) );
if (k1 < k2)
  return _test * (LiCrtRef - i_ex*(k1 - k2));
else
  return _test * LiCrtRef;
\end{lstlisting}
\notesRow{The conditional discards stripping current, modelling a one-way reaction.}
\end{ecmEntry}

\begin{ecmEntry}{CouplingMiec}
\inputsRow{coupled $\phi_{\text{Li}}$ (with gradient); $\sigma_{\text{ion}}$ (material);
$i_{\text{Li,ref}}$.}
\pdeRow{Imposes total flux $=$ applied + ionic ohmic contribution:
\[
  r = -\bigl(i_{\text{Li,ref}} + 10\,\sigma_{\text{ion}}\,\nabla\phi_{\text{Li}}\!\cdot\!\nvec\bigr)\psi.
\]}
\codeRow\begin{lstlisting}
return -_test * (LiCrtRef + 10 * sigma_ion * _grad_potLi * _normals);
\end{lstlisting}
\notesRow{Hard-coded $\times 10$ unit factor.}
\end{ecmEntry}

\begin{ecmEntry}{MixSEIntBV}
\inputsRow{$\phi_{\text{Li,electrode}}$ (material); coupled $\phi_e$;
$i_{\text{ex}},\,\alpha,\,F_{RT}=38.68$.}
\pdeRow{$\;\Delta\Phi = u + \phi_e - \phi_{\text{Li,electrode}};\;\; r = i_{\text{ex}}(k_1-k_2)\,\psi$.}
\codeRow\begin{lstlisting}
k1 = exp( alpha*F_RT*(u + potEn - LiPotEle) );
k2 = exp(-(1-alpha)*F_RT*(u + potEn - LiPotEle) );
return _test * i_ex * (k1 - k2);
\end{lstlisting}
\notesRow{$u$ is interpreted as the Li$^+$ potential.}
\end{ecmEntry}

\begin{ecmEntry}{MixSEMetalBV}
\inputsRow{coupled $\phi_{\text{ref}}$; $i_{\text{ex}},\alpha$;
$c_e$ (electron concentration material); $F_{RT}=38.68$.}
\pdeRow{$\;\Delta\Phi = \phi_{\text{ref}} + u;\;\; r = c_e\,i_{\text{ex}}(k_1-k_2)\,\psi$.}
\codeRow\begin{lstlisting}
k1 = exp( alpha*F_RT*(potRef + u) );
k2 = exp(-(1-alpha)*F_RT*(potRef + u) );
return _test * c_e * i_ex * (k1 - k2);
\end{lstlisting}
\notesRow{Pre-multiplies BV current by local electron concentration (mixed conductor).}
\end{ecmEntry}

\begin{ecmEntry}{MixSEVoidBV}
\inputsRow{Same as \textsc{MixSEMetalBV}.}
\pdeRow{Identical to \textsc{MixSEMetalBV} but rectified to $k_1>k_2$:
$\;r = c_e\,i_{\text{ex}}(k_1-k_2)\,\psi\,\text{ if }k_1>k_2,\,0\text{ otherwise.}$}
\codeRow\begin{lstlisting}
if (k1 > k2)
  return _test * c_e * i_ex * (k1 - k2);
else
  return 0;
\end{lstlisting}
\notesRow{Models deposition (forward only) into a void; no stripping.}
\end{ecmEntry}

\begin{ecmEntry}{MixSEMetalEn}
\inputsRow{coupled $\phi_{\text{metal}}$ (with gradient), $\phi_{\text{Li}}$;
$\sigma_m,c_e,i_{\text{ex}},\alpha,F_{RT}$.}
\pdeRow{$\;\Delta\Phi = \phi_{\text{Li}}+u;\;\;
r = \bigl(-10^4\sigma_m\nabla\phi_{\text{metal}}\!\cdot\!\nvec + c_e i_{\text{ex}}(k_1-k_2)\bigr)\psi$.}
\codeRow\begin{lstlisting}
return _test * (-10000*sigma_m * _grad_potMt * _normals
                + c_e * i_ex * (k1 - k2));
\end{lstlisting}
\notesRow{Combines BV with metal-side ohmic flux on the electron equation; $\times 10^4$
unit factor.}
\end{ecmEntry}

\begin{ecmEntry}{NeumannEn}
\inputsRow{$\phi_{\text{Li,electrode}}$ (material); coupled $\phi_{\text{Li}}$;
$i_{\text{inlet}},\,i_{\text{ex}},\,\alpha,\,F_{RT}$.}
\pdeRow{Cathode-side Neumann: applied current + BV at boundary.
$\;\Delta\Phi = \phi_{\text{Li}}+u-\phi_{\text{Li,electrode}};\;\;
r = (i_{\text{inlet}} + i_{\text{ex}}(k_1-k_2))\psi$.}
\codeRow\begin{lstlisting}
return _test * (i_inlet + i_ex * (k1 - k2));
\end{lstlisting}
\end{ecmEntry}

\begin{ecmEntry}{NeumannLi}
\inputsRow{coupled $\phi_e$ (with gradient); $\sigma_e$, $i_{\text{inlet}}$ (AD materials).}
\pdeRow{$\;r = (10^4\,\sigma_e\,\nabla u\!\cdot\!\nvec - i_{\text{inlet}})\psi$.}
\codeRow\begin{lstlisting}
return _test * (10000 * sigma_e * _grad_u * _normals - i_inlet);
\end{lstlisting}
\notesRow{Pairs with \textsc{NeumannEn}; subtracts the imposed inlet current and adds
electron ohmic flux. $\times 10^4$ unit factor.}
\end{ecmEntry}

\begin{ecmEntry}{SingleSEElectrodeBV}
\inputsRow{$\phi_{\text{Li,electrode}}$ (material); $i_{\text{ex}},\alpha,F_{RT}$.}
\pdeRow{$\;\Delta\Phi = u-\phi_{\text{Li,electrode}};\;\; r = i_{\text{ex}}(k_1-k_2)\psi$.}
\codeRow\begin{lstlisting}
return _test * i_ex * (k1 - k2);
\end{lstlisting}
\notesRow{Single-conductor SE/electrode interface; $u$ is the full overpotential variable.}
\end{ecmEntry}

\begin{ecmEntry}{SingleSEIntLiBV}
\inputsRow{coupled $\phi_{\text{Li}}$; $i_{\text{ex}},\alpha,F_{RT}$.}
\pdeRow{$\;\Delta\Phi = \phi_{\text{Li}}-u;\;\; r = i_{\text{ex}}(k_1-k_2)\psi$.}
\codeRow\begin{lstlisting}
return _test * i_ex * (k1 - k2);
\end{lstlisting}
\end{ecmEntry}

\begin{ecmEntry}{SingleSEIntSEBV}
\inputsRow{coupled $\phi_e$; $i_{\text{ex}},\alpha,F_{RT}$.}
\pdeRow{Mirror of the previous entry with sign flipped:
$\Delta\Phi = u-\phi_e;\;\;r=i_{\text{ex}}(k_1-k_2)\psi$.}
\codeRow\begin{lstlisting}
return _test * i_ex * (k1 - k2);
\end{lstlisting}
\end{ecmEntry}

%==========================================================================
\section{Interface Kernels (\texttt{src/interfacekernels/})}
%==========================================================================

Interface kernels return one residual on the \emph{element} side and one on the
\emph{neighbor} side; we list both.

\begin{ecmEntry}{ADInterfaceDiffusion}
\inputsRow{$D,\,D_{\text{neighbor}}$ (AD materials).}
\pdeRow{Enforces Fickian-flux continuity across the interface:
\[
  r_{\text{elem}} = -D_{\text{neighbor}}\nabla u_{\text{neighbor}}\cdot\nvec\,\psi,\quad
  r_{\text{nbr}}  = +D\,\nabla u\cdot\nvec\,\psi_{\text{nbr}}.
\]}
\codeRow\begin{lstlisting}
ELEMENT: -_test * D_neighbor * _grad_neighbor_value * _normals;
NEIGHBOR: +_test_neighbor * D * _grad_u * _normals;
\end{lstlisting}
\notesRow{Choice of signs makes the two halves cancel when $D\nabla u = D_{\text{n}}\nabla u_{\text{n}}$.}
\end{ecmEntry}

\begin{ecmEntry}{PenaltyIntDiffusion}
\inputsRow{\texttt{penalty} $\lambda_p$.}
\pdeRow{Nitsche-style penalty: $\;r_{\text{elem}} = +\lambda_p(u-u_n)\psi,\;\;
r_{\text{nbr}} = -\lambda_p(u-u_n)\psi_{\text{nbr}}$.}
\codeRow\begin{lstlisting}
ELEMENT:  _test * penalty * (u - u_neighbor);
NEIGHBOR: -_test_neighbor * penalty * (u - u_neighbor);
\end{lstlisting}
\notesRow{Used to glue two scalar variables across an interface.}
\end{ecmEntry}

\begin{ecmEntry}{IonicSEInterface}
\inputsRow{$\sigma_{\text{ion}},\sigma_{\text{metal}},i_{\text{ex}},\alpha$ (AD materials);
$F_{RT}=38.68$.}
\pdeRow{Asymmetric BV/ohmic split:
$\Delta\Phi = u + u_n;$
\[
 r_{\text{elem}} = +10^4\,\sigma_{\text{metal}}\nabla u_n\!\cdot\!\nvec\,\psi,\quad
 r_{\text{nbr}}  = +i_{\text{ex}}(k_1-k_2)\,\psi_{\text{nbr}}.
\]}
\codeRow\begin{lstlisting}
ELEMENT:  _test*10000*sigma_metal*_grad_neighbor_value*_normals;
NEIGHBOR: _test_neighbor * i_ex * (k1 - k2);
\end{lstlisting}
\notesRow{$\sigma_{\text{ion}}$ is declared but unused in the residual.}
\end{ecmEntry}

\begin{ecmEntry}{MiecSEInterfaceEn}
\inputsRow{coupled $\phi_{\text{Li}}$; $c_e,\sigma_{\text{ion}},\sigma_e,\sigma_m,
i_{\text{ex}},\alpha,F_{RT}$.}
\pdeRow{Full MIEC interface for the electron equation (BV on both sides + ohmic flux on each):
$\Delta\Phi = \phi_{\text{Li}} + u$,
\[
\begin{aligned}
 r_{\text{elem}} &= \bigl(c_e i_{\text{ex}}(k_1-k_2) - 10^4\sigma_m \nabla u_n\!\cdot\!\nvec\bigr)\psi,\\
 r_{\text{nbr}}  &= \bigl(c_e i_{\text{ex}}(k_1-k_2) + 10^4\sigma_e \nabla u\!\cdot\!\nvec\bigr)\psi_{\text{nbr}}.
\end{aligned}
\]}
\codeRow\begin{lstlisting}
ELEMENT: _test*(c_e*i_ex*(k1-k2)
               - 10000*sigma_metal*_grad_neighbor_value*_normals);
NEIGHBOR: _test_neighbor*(c_e*i_ex*(k1-k2)
               + 10000*sigma_e*_grad_u*_normals);
\end{lstlisting}
\end{ecmEntry}

\begin{ecmEntry}{MiecSEInterfacePrs}
\inputsRow{Same as \textsc{MiecSEInterfaceEn} plus mechanical pressure $P$ (Real) and
$V_m/RT = 0.0052$.}
\pdeRow{Adds a mechanical contribution to $\Delta\Phi$:
$\Delta\Phi = F_{RT}(\phi_{\text{Li}}+u) - (V_m/\Rg T)\,P$.
Otherwise identical residuals to \textsc{MiecSEInterfaceEn}.}
\codeRow\begin{lstlisting}
k1 = exp( alpha*( F_RT*(potLi+u) - Vm_RT*P ) );
k2 = exp(-(1-alpha)*( F_RT*(potLi+u) - Vm_RT*P ) );
\end{lstlisting}
\notesRow{Couples electrochemistry to mechanical pressure (electrochemo-mechanical
interface). The factor $V_m/\Rg T$ in mV$^{-1}\!\cdot$\,Pa$^{-1}$ requires consistent units.}
\end{ecmEntry}

\begin{ecmEntry}{CZMDiffusiveVariableKernelBase / SmallStrain}
\inputsRow{component, displacements (vector coupled), \texttt{base\_name}, \texttt{flux\_global}
material, \texttt{include\_gap}; reads materials \texttt{flux\_global},
\texttt{dflux\_dvariablejump\_global}, \texttt{dflux\_djump\_global}.}
\pdeRow{Imposes a generic interface flux:
\[
  r_{\text{elem}} = -j_{\text{global}}\,\psi,\quad r_{\text{nbr}}=+j_{\text{global}}\,\psi_{\text{nbr}}.
\]}
\codeRow\begin{lstlisting}
ELEMENT: -_test * flux_global;
NEIGHBOR: +_test_neighbor * flux_global;
\end{lstlisting}
\notesRow{The base provides hand-coded \texttt{computeDResidualDVariable} and
\texttt{computeDResidualDDisplacement} hooks; the small-strain subclass implements them.
Used together with the cohesive-zone material stack.}
\end{ecmEntry}

%==========================================================================
\section{Mortar Constraints (\texttt{src/constraints/})}
%==========================================================================

All five are \texttt{ADMortarConstraint}s with three residual rows: \emph{primary},
\emph{secondary}, and \emph{lower-dimensional} (the Lagrange multiplier $\lambda$).

\begin{ecmEntry}{EqualNormalFluxConstraint}
\inputsRow{primary/secondary diffusivity material name (default \texttt{diffusivity});
\texttt{primary\_tensor}, \texttt{secondary\_tensor} bool flags.}
\pdeRow{Enforces $\bigl[\!\!\bigl[D\nabla u\bigr]\!\!\bigr]\!\cdot\!\nvec = 0$ via Lagrange
multiplier $\lambda$:
\[
  r_p = \lambda\psi_p,\;\; r_s = \lambda\psi_s,\;\;
  r_\lambda = \bigl((D_p\nabla u_p - D_s\nabla u_s)\!\cdot\!\nvec + \lambda\bigr)\psi.
\]}
\notesRow{Supports scalar or tensor $D$ on each side independently.}
\end{ecmEntry}

\begin{ecmEntry}{GapDisplacementConductanceConstraint}
\inputsRow{$k$ or $k(\text{gap})$; \texttt{conductanceType}\,$\in$\,
\{CONDUCTANCE, RCT, EXCHANGE\_CURRENT\_DENSITY\};
$\F=96485.33,\Rg=8.3145,T=298$;
displacements; concentration $c$; flags
\texttt{include\_equilibrium\_potential}, \texttt{include\_gap}, \texttt{include\_concentration};
\texttt{surfaceType}\,$\in$\,\{PRIMARY, SECONDARY\};
\texttt{computeType}\,$\in$\,\{LINEAR\_BUTLER\_VOLMER, BUTLER\_VOLMER\}.}
\pdeRow{The heart of the displaced-gap electrochemical contact. Defining
$\Delta\Phi = \sgn(\,\cdot\,)(u_s - u_p) - U_{eq}$ and $\mathrm{RTF}=\Rg T/\F$:
\[
\begin{aligned}
\text{linear:}\quad& r_\lambda = \bigl(\lambda + \tfrac{i_0}{\mathrm{RTF}}\,\Delta\Phi\bigr)\psi,\\
\text{full BV:}\quad& r_\lambda = \bigl(\lambda + 2 i_0\,\sinh(\Delta\Phi/(2\,\mathrm{RTF}))\bigr)\psi.
\end{aligned}
\]
$r_p = \lambda\psi_p,\; r_s = -\lambda\psi_s$. If $c<0$ the constraint clamps to $\lambda$
only (no reaction).}
\notesRow{Supports gap-modulated current via $k(\text{gap})$ using dual derivatives at
phys-points. Three different input conventions converge on a single $i_0$.}
\end{ecmEntry}

\begin{ecmEntry}{GapEquilibriateConstraint}
\inputsRow{$k$ or $k(\text{gap})$; \texttt{primary\_mat\_prop}, \texttt{secondary\_mat\_prop}
(default \texttt{equilibrium\_potential}); \texttt{one\_sided}\,$\in$\,
\{\textsc{none}, \textsc{primary$\to$secondary}, \textsc{secondary$\to$primary}\};
\texttt{prefactor}\,$\in$\,\{\textsc{rt/f}, \textsc{rt}, \textsc{none}\}.}
\pdeRow{Equilibrates a chosen material property across the gap:
\[
  r_\lambda = \bigl(\lambda + \tfrac{k}{\text{prefactor}}\,(p_p - p_s)\bigr)\psi,
\]
with optional one-sided ratchet (only acts in one direction).}
\end{ecmEntry}

\begin{ecmEntry}{PressureLMConductanceConstraint}
\inputsRow{$k$, threshold pressure $p_{\text{th}}=10^{-15}$, contact pressure (lower-dim
nonlinear var, the contact LM).}
\pdeRow{Gates a linear conductance by the contact LM:
\[
  r_\lambda = \begin{cases}\bigl(k(u_p-u_s)-\lambda\bigr)\psi, & p_{\text{contact}}>p_{\text{th}},\\
   -\lambda\psi, & \text{otherwise.}\end{cases}
\]}
\notesRow{No contact $\Rightarrow$ no electrical/thermal conductance.}
\end{ecmEntry}

\begin{ecmEntry}{ScaledBCConstraint}
\inputsRow{\texttt{scale}, \texttt{primary} (bool).}
\pdeRow{Scales an existing mortar $\lambda$ to a different equation:
$r_p = -\lambda\,$\texttt{scale}$\,\psi_p$ (if \texttt{primary}) or
$r_s = -\lambda\,$\texttt{scale}$\,\psi_s$ (otherwise). Sets
\texttt{compute\_lm\_residuals=false}.}
\notesRow{Used to couple, e.g., a stress LM into the chemistry equation.}
\end{ecmEntry}

%==========================================================================
\section{Materials (\texttt{src/materials/})}
%==========================================================================

\subsection*{1. Butler--Volmer kinetics and equilibrium potential}

\begin{ecmEntry}{ADButlerVolmerMaterial}
\label{bv-mat}
\inputsRow{coupled $\phi_{\text{electrolyte}},\phi_{\text{electrode}}$;
$\F=96.485\;\mathrm{C/mmol},\Rg=8.3145,T=298$; \texttt{include\_equilibrium} bool;
$i_0$ (material); surface-to-volume ratio $a_v$; volume fraction $\eps_v$; \texttt{specific\_area}.}
\pdeRow{Produces the BV kinetic quantities consumed by \textsc{ADButlerVolmerForce}:
\[
\begin{aligned}
\eta &= \phi_{\text{electrode}} - \phi_{\text{electrolyte}} - U_{eq},\\
i_{\text{BV}} &= 2\,i_0\,\sinh(\eta\F/(2\Rg T)),\\
J_{\text{BV}} &= i_{\text{BV}}/\F,\\
J_{\text{BV}}^{\text{vol}} &= a_v\,J_{\text{BV}},\\
i_{\text{BV}}^{\text{vol}} &= a_v\,\eps_v\,i_{\text{BV}}.
\end{aligned}
\]}
\notesRow{Faraday default in C/mmol (instead of C/mol); intended for input files using mmol/cm$^3$.}
\end{ecmEntry}

\begin{ecmEntry}{ADComputeEquilibriumPotential}
\inputsRow{coupled $c$; \texttt{reaction\_rate} (Real or Function $f(\bar c)$);
\texttt{is\_ocv} bool; \texttt{potential}\,$\in$\,\{\textsc{dehoff}, \textsc{bower}\};
flags \texttt{include\_reaction\_rate}, \texttt{include\_conc},
\texttt{include\_mechanical\_effects},
\texttt{exclude\_elastic\_contribution}; \texttt{nucleation\_overpotential}.}
\pdeRow{Outputs three properties: state of charge $\bar c$ (clipped to $[10^{-3},1-10^{-3}]$),
chemical potential, and equilibrium potential. Defining $\bar c = c/c_0$:
\[
\begin{aligned}
\mu_{\text{chem}} &= \begin{cases}
   \Rg T \ln(c_0/(c_0-c)), & \text{DEHOFF},\\
   \Rg T \ln \bar c,        & \text{BOWER},
 \end{cases}\\[4pt]
\mu_{\text{chem}} &\mathrel{+}= \mu_{\text{swell}} + \mu_{\text{elastic}}\;
    [\text{if mech. effects on}],\\[2pt]
U_{eq} &= \begin{cases}f(\bar c),& \text{is\_ocv},\\ f(\bar c)\,\Rg T/\F,& \text{else,}\end{cases}\\
U_{eq} &\mathrel{-}= \mu_{\text{chem}}/\F.
\end{aligned}
\]}
\notesRow{Covers OCV-table and Nernstian forms with optional mechanical coupling. Two
chemical-potential conventions (DEHOFF vs BOWER) differ in their treatment of the
configurational entropy.}
\end{ecmEntry}

\begin{ecmEntry}{ADExchangeCurrentDensityMaterial}
\inputsRow{coupled $c$; type\,$\in$\,\{CONSTANT, LI\_INSERTION\}; $i_{\text{ref}},\,c_{\max}$.}
\pdeRow{Standard intercalation form:
\[
  i_0 = \begin{cases} i_{\text{ref}}, & \text{CONSTANT},\\
   2\,i_{\text{ref}}\sqrt{\bar c\,(1-\bar c)}, & \text{LI\_INSERTION,}\end{cases}
\]
with $\bar c=c/c_{\max}$ clipped to $[10^{-3},1-10^{-3}]$.}
\end{ecmEntry}

\begin{ecmEntry}{ADParticleSize}
\inputsRow{$R$ (Real or function); particleType\,$\in$\,\{\textsc{spherical},\textsc{cylindrical}\};
$\eps_v$.}
\pdeRow{$\;a_v = \frac{3}{R}\;\text{(sphere)}\quad\text{or}\quad \frac{2}{R}\;\text{(cylinder)};\quad
   \texttt{volume\_fraction}=\eps_v.$}
\notesRow{Provides $a_v$ for the BV-volume source.}
\end{ecmEntry}

\subsection*{2. Hyperelastic / visco-plastic / swelling stack}

\begin{ecmEntry}{ADIsoTropicHyperViscoSwellingBase (base)}
\label{mat-swelling-base}
\inputsRow{\texttt{isotropic\_swelling} bool; \texttt{rate\_form} bool;
coupled $c$, reference $c_{\text{ref}}=0$;
$\bm{\alpha}\!\in\!\mathbb{R}^3$ (sums to 1); molar volume $\Omega$;
boundary names; \texttt{use\_deformed\_directions} bool.}
\pdeRow{Multiplicative split of the deformation gradient:
$\mathbf{F} = \mathbf{F}^e\mathbf{F}^p\mathbf{F}^g$. Growth Jacobian
$J_g = 1 + \Omega(c-c_{\text{ref}})$, with anisotropic $\mathbf{F}^g$
\[
  \mathbf{F}^g = \alpha_n J_g\,\nvec\!\otimes\!\nvec + \alpha_t J_g\,\bm{t}\!\otimes\!\bm{t} +
                \alpha_z J_g\,\hat{\bm z}\!\otimes\!\hat{\bm z}\quad\text{or}\quad
  J_g^{1/3}\mathbf{I}\;[\text{isotropic}].
\]
Logarithmic elastic kinematics: $\bm{F}_i=\bm{F}^p\bm{F}^g$, trial Mandel stress
$\bm{M}_e^{\text{tr}} = 2\mu\,\mathrm{dev}\,\bm{E}^e + \kappa\,\ln(\det\bm{F}/J_g)\bm{I}$.
The base produces:
\[
\begin{aligned}
\mu_{\text{swell}} &= -\Omega\,\bm{\sigma}_g:\bm{F}^g,\\
\mu_{\text{elastic}} &= \tfrac{1}{2}\Omega\,\bm{E}^e:\mathbb{C}:\bm{E}^e.
\end{aligned}
\]
These chemical potentials are consumed by \textsc{ADComputeEquilibriumPotential} and
\textsc{ADStressBasedChemicalPotential}.}
\notesRow{Anisotropic split uses a swell normal computed by closest-edge construction at
specified boundaries.}
\end{ecmEntry}

\begin{ecmEntry}{ADIstropicElasticSwelling}
\inputsRow{Inherits base.}
\pdeRow{Pure elastic + swelling kinematics: returns
$\bm{\sigma} = \bm{R}\,\bm{M}_e^{\tau}\,\bm{R}^{\!\top}\,J_g/\det\bm{F}$ with
$\bm{M}_e^\tau = \mathbb{C}:\bm{E}^e$. Same $\mu_{\text{swell}}$, $\mu_{\text{elastic}}$
as base.}
\end{ecmEntry}

\begin{ecmEntry}{ADIsoTropicHyperVisco}
\inputsRow{rate exponent $m$, reference rate $\dot\gamma_{\text{ref}}$,
$Y_0,\,H_0,\,Y_{\text{sat}}$, hardening exponent $a$.}
\pdeRow{Weber--Anand 1990 rate-dependent isotropic visco-plastic model. Saturation
hardening $\dot Y = H_0(1-Y/Y_{\text{sat}})^a\,\dot\gamma$. Return-mapping residual:
\[
  r_{\text{rm}} = \frac{\sigma_{e,\text{tr}} - \mathrm{d}t\,\dot\gamma\cdot 3\mu - Y(\dot\gamma)
    \,(\dot\gamma/\dot\gamma_{\text{ref}})^m}{3\mu/3}.
\]
Updates: $\bm{F}^p = \exp(\Delta\bm{p})\bm{F}^p_{\text{old}}$,
$\bm{\sigma} = \bm{R}\,\bm{M}_e^\tau\,\bm{R}^{\!\top}/\det\bm{F}$.}
\notesRow{Produces \texttt{plastic\_strain}, \texttt{Fp}, mandel stress, log-elastic strain.}
\end{ecmEntry}

\begin{ecmEntry}{ADIsoTropicHyperViscoCreep}
\pdeRow{Same skeleton as the previous entry, with Arrhenius-activated reference rate
$\dot\gamma_0 = A\,\exp(-Q/\Rg T)$ and saturation strength dependent on rate
$Y_{\text{sat},\dot\gamma} = Y_{\text{sat}}(\dot\gamma/\dot\gamma_{\text{ref}})^n$.}
\notesRow{Thermally-activated viscoplasticity.}
\end{ecmEntry}

\begin{ecmEntry}{ADIsoTropicHyperViscoSwelling}
\pdeRow{Combines \textsc{ADIsoTropicHyperVisco} return-mapping with the swelling
kinematics base ($\bm{F}=\bm{F}^e\bm{F}^p\bm{F}^g$). Room-temperature swelling +
viscoplasticity.}
\end{ecmEntry}

\begin{ecmEntry}{ADIsoTropicHyperViscoCreepSwelling}
\pdeRow{Combines Arrhenius-activated visco-plasticity with the swelling base.}
\notesRow{Most general member of the stack.}
\end{ecmEntry}

\begin{ecmEntry}{ADIsoTropicHyperViscoBase}
\pdeRow{Internal base: polar decomposition $\bm{F}^e = \bm{F}\bm{F}^{p\,-1}$, log strain
$\bm{E}^e=\sum\ln(\lambda_i)\bm{N}_i\bm{N}_i^{\!\top}$, trial Mandel stress
$\bm{M}_e^{\text{tr}} = \mathbb{C}\bm{E}^e$, return mapping, and exp-update of $\bm{F}^p$.}
\notesRow{Core kinematic engine reused by all four visco-plastic variants.}
\end{ecmEntry}

\subsection*{3. Diffusivity and conductivity providers}

\begin{ecmEntry}{IsotropicDiffusionMaterial / PorousIsoTropicDiffusionMaterial}
\inputsRow{\texttt{diffusion\_coef} $D_0$; subclass adds $\eps_v,\,\beta_{\text{Brug}}=1.5,\,\tau$;
\texttt{diffusion\_model}\,$\in$\,\{\textsc{bruggeman}, \textsc{tortuosity}, \textsc{user\_defined}\}.}
\pdeRow{$D_{\text{eff}} = D_0\cdot c(\eps_v)$, with
\[
  c(\eps_v) = \begin{cases}
    \eps_v^{\beta_{\text{Brug}}}, & \text{Bruggeman},\\
    \eps_v/\tau, & \text{tortuosity},\\
    c_{\text{user}}, & \text{user-defined.}
  \end{cases}
\]}
\end{ecmEntry}

\begin{ecmEntry}{AnisotropicDiffusivity}
\inputsRow{$k_0,\,w,\,\lambda$; coupled gradient variable $v$.}
\pdeRow{Kobayashi-style angle-dependent diffusivity:
$\;\theta = \arctan(\partial_y v / \partial_x v),\;\;
D = k_0(1 + w\cos(\lambda\theta))^2$.}
\notesRow{Returns scalar $D$, not a tensor.}
\end{ecmEntry}

\begin{ecmEntry}{DiffusionAlongPrincipalDirectionsMaterial}
\inputsRow{$D_1,D_2,D_3$ (vector components); reads \texttt{swell\_normal} from base.}
\pdeRow{Builds an anisotropic diffusivity tensor from three eigen-directions tied to the
swelling normal:
\[
  \mathbf{D} = D_1\,\bm{m}_1\bm{m}_1^{\!\top} + D_2\,\bm{m}_2\bm{m}_2^{\!\top} + D_3\,\bm{m}_3\bm{m}_3^{\!\top},
\]
with $\bm{m}_1 = \texttt{swell\_normal}$, $\bm{m}_3 = (0,0,1)$, $\bm{m}_2 = \bm{m}_1\!\times\!\bm{m}_3$.}
\end{ecmEntry}

\begin{ecmEntry}{AnisotropyMaterialConst}
\inputsRow{Nine components $\sigma_{ij}$.}
\pdeRow{Assembles a constant rank-2 tensor as an AD material property.}
\notesRow{Declared AD but populated from non-AD constants; sensitivities are trivial.}
\end{ecmEntry}

\begin{ecmEntry}{RotatedConductivityTensor}
\inputsRow{Nine components $\sigma_{ij}$, three Euler angles.}
\pdeRow{Builds $\bm{\sigma}$, then computes $\bm{\sigma}_{\text{rot}} = \bm{R}^{\!\top}\bm{\sigma}\bm{R}$
and exports the 9 scalar components as separate AD properties.}
\end{ecmEntry}

\begin{ecmEntry}{Ionics}
\inputsRow{$\sigma_{\text{ion}},\sigma_{gb}$.}
\pdeRow{Constant SE material: writes $\sigma_{\text{ion}},\sigma_{gb}$.}
\end{ecmEntry}

\begin{ecmEntry}{Miec}
\inputsRow{$\sigma_{\text{ion}},\sigma_e,\sigma_{gb}$.}
\pdeRow{Constant MIEC material.}
\end{ecmEntry}

\begin{ecmEntry}{MiecSEAgC / MixAgC}
\inputsRow{Ten Reals: $\sigma_{\text{ion,SE}},\sigma_{\text{ion,AgC}},\sigma_{e,\text{SE}},\sigma_{e,\text{AgC}},
i_{\text{inlet}},i_{\text{ex}},\alpha,c_e,\phi_{\text{Li,anode}},\phi_{\text{Li,cathode}}$.}
\pdeRow{Bulk-property bag for SE/Ag-C composite domain: writes all ten as AD properties.}
\notesRow{\textsc{MixAgC} is byte-for-byte equivalent to \textsc{MiecSEAgC} (only the
class registration name differs); pick one and remove the duplicate.}
\end{ecmEntry}

\begin{ecmEntry}{MixSE}
\inputsRow{Nine Reals: $\sigma_{\text{ion}},\sigma_m,\sigma_e,i_{\text{inlet}},i_{\text{ex}},\alpha,c_e,
\phi_{\text{Li,anode}},\phi_{\text{Li,cathode}}$.}
\pdeRow{Bulk material for SE/metal/void multi-conductor problem.}
\end{ecmEntry}

\begin{ecmEntry}{SingleSE}
\inputsRow{$\sigma_{\text{ion}},\sigma_{gb},\sigma_m,i_{\text{applied}},i_{\text{ex}},\alpha,
\phi_{\text{Li,anode}},\phi_{\text{Li,cathode}}$.}
\pdeRow{Single-conductor SE variant.}
\end{ecmEntry}

\begin{ecmEntry}{PaperHongli}
\inputsRow{coupled $c_{\text{Li}}$; constants $c_1,c_2$.}
\pdeRow{Empirical concentration-dependent diffusivity from a literature source:
\[
  D = c_1\cdot 10^{6 c_2 c_{\text{Li}}},\qquad \sigma_{\text{ion}} = 0.11\,D.
\]}
\notesRow{Magic factors $6$ and $0.11$ inherited from the cited paper; verify before reuse.}
\end{ecmEntry}

\subsection*{4. Cohesive zone model}

\begin{ecmEntry}{CZMcomputeLocalFluxBase / CZMComputeVariableJumpBase (bases)}
\pdeRow{Define the contract for cohesive-zone diffusive variables: \emph{LocalFluxBase}
declares \texttt{interface\_flux}, $\partial$flux/$\partial$jump, $\partial$flux/$\partial$displacement;
\emph{VariableJumpBase} declares total rotation, global variable jump, and the local
interface variable jump. Derived classes provide
\texttt{computeInterfaceFluxAndDerivatives} / \texttt{computeLocalVariableJump}.}
\end{ecmEntry}

\begin{ecmEntry}{CZMComputeVariableJumpSmallStrain}
\pdeRow{Trivial small-strain implementation:
\[
  [\![u]\!]^{\text{local}} = [\![u]\!]^{\text{global}} = u_{\text{neighbor}} - u_{\text{master}}.
\]}
\notesRow{Registered to \texttt{TensorMechanicsApp} (not \texttt{ecmApp}).}
\end{ecmEntry}

\begin{ecmEntry}{CZMcomputeGlobalFluxBase / CZMComputeGlobalFluxSmallStrain}
\pdeRow{Rotates the local-frame flux to the global frame:
\[
  j_{\text{global}} = j_{\text{interface}}, \quad
  \frac{\partial j_{\text{global}}}{\partial[\![u]\!]} = \frac{\partial j_{\text{interface}}}{\partial[\![u]\!]},\quad
  \frac{\partial j_{\text{global}}}{\partial[\![\bm{u}]\!]} = \bm{R}_{\text{tot}}\,\frac{\partial j_{\text{interface}}}{\partial[\![\bm{u}]\!]}.
\]}
\notesRow{Small-strain version is a near-identity pass-through.}
\end{ecmEntry}

\begin{ecmEntry}{CZMButlerVolmer}
\inputsRow{$d$ (max separation), $k(\text{jump})$ function, conductanceType,
\texttt{computeType}\,$\in$\,\{LBV, BV, LITHIUM\_INSERTION\}, $\F,\Rg,T$, $c_{\text{ref}}=1$,
$U_{eq}$ from \texttt{reaction\_rate\_function}, surfaceType.}
\pdeRow{Cohesive-zone Butler--Volmer with gap-modulated current. Defining
$\mathrm{RTF}=\Rg T/\F$, $\bar c = c_{\text{coupled}}/c_{\text{ref}}$,
$\sgn = +1$ on secondary, $-1$ on primary, $\Delta\Phi=\sgn\cdot[\![u]\!]_v - U_{eq}$:
\[
\begin{aligned}
\text{linear:} && j &= i_0\,\Delta\Phi/\mathrm{RTF},\\
\text{full BV:} && j &= 2\,i_0\,\sinh(\Delta\Phi/(2\,\mathrm{RTF})).
\end{aligned}
\]
If \texttt{include\_gap} and the normal jump $[\![u]\!]_n < d$:
$j \leftarrow j\cdot(d - [\![u]\!]_n)$ (linear taper to zero at $d$).}
\notesRow{The \emph{conductanceType} switch maps three input conventions
(conductance, $R_{ct}$, $i_0$) onto one $i_0$ used in the formula above.}
\end{ecmEntry}

\begin{ecmEntry}{SalehaniIrani3DCTractionViscosity}
\pdeRow{Salehani--Irani 3DC cohesive traction with viscous regularisation. Augments the
base law $T_i$ by a rate term $T_i \mathrel{*\!=}\bigl(1 + \tfrac{\mu_i}{T_{\max,i}}\,\dot{[\![u]\!]_i}\bigr)$.}
\end{ecmEntry}

\begin{ecmEntry}{GaoBower3DCTraction}
\inputsRow{normal/tangential viscosity, $r=0,\,q=1$.}
\pdeRow{Gao--Bower 3D cohesive law (Xu--Needleman style) with
$\bar n_i = [\![u]\!]_i/[\![u]\!]_{0,i}$:
\[
\begin{aligned}
T_0 &= T_{\max}\,e^{1-\bar n_0}\Bigl[\bar n_0 e^{-\sum_{i>0}\bar n_i^2} +
       \tfrac{1-q}{r-1}\bigl(1 - e^{-\sum_{i>0}\bar n_i^2}\bigr)(r - \bar n_0)\Bigr],\\
T_i &= T_{\max}\cdot 2\,([\![u]\!]_{0,0}/[\![u]\!]_{0,1})\,\bar n_i\Bigl(q + \tfrac{r-q}{r-1}\bar n_0\Bigr)
       e^{-\sum_{j>0}\bar n_j^2}\,e^{1-\bar n_0}.
\end{aligned}
\]
Then a viscous rate factor $T_i\mathrel{*\!=}(1 + \mu_i\,\dot{[\![u]\!]_i}/T_{\max})$ is applied.}
\end{ecmEntry}

%==========================================================================
\section{Aux Kernels (\texttt{src/auxkernels/})}
%==========================================================================

\begin{ecmEntry}{ADDiffusionFluxAux}
\inputsRow{component $i$, coupled $u$, $D$ (AD).}
\pdeRow{$\;j_i = -D\,\partial_i u$ (scalar diffusivity component).}
\notesRow{\texttt{raw\_value(D)} drops AD sensitivity.}
\end{ecmEntry}

\begin{ecmEntry}{AnisoTropicDiffusionFlux}
\inputsRow{component, coupled $u$, $\bm{D}$ (AD tensor).}
\pdeRow{$\;j_i = -(\bm{D}\nabla u)_i$.}
\end{ecmEntry}

\begin{ecmEntry}{DiffusionFluxNormalToBoundaryAux}
\inputsRow{coupled $u$, $D$ (AD).}
\pdeRow{$\;j_n = -D\,\nabla u\!\cdot\!\nvec$ on a boundary auxiliary.}
\end{ecmEntry}

\begin{ecmEntry}{AnisoDiffusionFluxNoramltoBoundary}
\inputsRow{coupled $u$, $\bm{D}$ (AD tensor).}
\pdeRow{$\;j_n = -(\bm{D}\nabla u)\!\cdot\!\nvec$.}
\notesRow{Filename retains the original misspelling (``Noraml'').}
\end{ecmEntry}

\begin{ecmEntry}{ADRankTwoTensorNormalToBoundary}
\inputsRow{tensor $\bm{T}$ (AD); reads boundary normals.}
\pdeRow{Intent: compute $\nvec^{\!\top}\bm{T}\nvec$.}
\codeRow\begin{lstlisting}
temp_i = n . T.column(i);
p1 = (raw(temp1), raw(temp2), raw(temp3));   // <-- comma operator
p2 = (raw n);
return p1 * p2;
\end{lstlisting}
\notesRow{\textbf{Likely bug}: \texttt{p1} is constructed using the comma operator, which
produces only the last value, not a 3-vector. The dot product is therefore not the intended
$\nvec^{\!\top}\bm{T}\nvec$. Should be replaced by \texttt{Point(raw(temp1), raw(temp2), raw(temp3))}.}
\end{ecmEntry}

\begin{ecmEntry}{ComputeSurfaceConcentrationAux}
\inputsRow{coupled $\bar c$; flux $j$ (AD); $D,\,R,\,$\texttt{scale}.}
\pdeRow{$\;c_{\text{surf}} = \bar c + R\,j/(\text{scale}\,D)$
(2-term polynomial recovery from average + flux).}
\end{ecmEntry}

\begin{ecmEntry}{CurrentDensity}
\inputsRow{component, conductivity $\sigma$ (AD scalar), coupled $\phi$.}
\pdeRow{$\;i_i = -\sigma\,\partial_i\phi\cdot 10^4$;\, sign flipped if material name is
\texttt{ionic\_conductivity}.}
\notesRow{Hard-coded $10^4$ unit conversion + name-based sign convention is fragile but
intentional.}
\end{ecmEntry}

\begin{ecmEntry}{CurrentDensityAnisotropy}
\inputsRow{component, $\bm{\sigma}$ (AD tensor), coupled $\phi$.}
\pdeRow{$\;\bm{i} = -10\,\bm{\sigma}\nabla\phi$; sign flipped if material name is
\texttt{ionic\_conductivity}.}
\notesRow{Different prefactor ($\times 10$) than \textsc{CurrentDensity} ($\times 10^4$);
verify unit consistency.}
\end{ecmEntry}

\begin{ecmEntry}{LiPotInSE / LiPotMiec}
\inputsRow{coupled $\phi_{\text{Li}},\phi_e$.}
\pdeRow{$\;u_{\text{aux}} = \phi_{\text{Li}} + \phi_e$ (total Li chemical potential).}
\notesRow{The two are byte-identical; the only practical difference is intended block
restriction.}
\end{ecmEntry}

%==========================================================================
\section{Postprocessors (\texttt{src/postprocessors/})}
%==========================================================================

\begin{ecmEntry}{EnergyRatePostprocessor}
\inputsRow{another postprocessor $p$, time-step $\Delta t$.}
\pdeRow{Returns $\bigl|(|p| - |p_{\text{old}}|)/|p|\bigr|/\Delta t$, an order-of-magnitude
rate-of-change diagnostic.}
\notesRow{Often used as a transient stopping criterion.}
\end{ecmEntry}

\begin{ecmEntry}{FBulk}
\inputsRow{coupled order parameter $u$; $A$ (material); \texttt{len\_scale}=1.}
\pdeRow{$\;\int_\Omega -A\,u^2(1-u)^2\,\dV$ (negated double-well bulk energy).}
\notesRow{Sign convention is unusual; verify against the energy-balance setup.}
\end{ecmEntry}

\begin{ecmEntry}{SideTotCrnt}
\inputsRow{conductivity $\sigma$ (AD).}
\pdeRow{$\;\int_{\partial\Omega} 100\,\sigma\,\nabla u\!\cdot\!\nvec\,\dS$;
sign flipped if material name is \texttt{ionic\_conductivity}.}
\notesRow{$\times 100$ unit factor (yet another convention vs.\ $\times 10$ and $\times 10^4$
in the related aux kernels).}
\end{ecmEntry}

%==========================================================================
\section{MultiApps (\texttt{src/multiapps/})}
%==========================================================================

\begin{ecmEntry}{QuadratureMultiApp}
\inputsRow{block restriction (suppresses manual \texttt{positions} / \texttt{positions\_file}).}
\pdeRow{Auto-generates one sub-app per quadrature point of the master mesh on the selected
blocks: iterates over local active elements, builds the default quadrature, pushes each
$\bm{x}_q$ into the position list, then \texttt{allgather}+sort.}
\notesRow{Used for sub-grid (e.g.\ particle-scale) coupling.}
\end{ecmEntry}

%==========================================================================
\section*{Cross-cutting observations}
\addcontentsline{toc}{section}{Cross-cutting observations}
%==========================================================================

\begin{enumerate}
  \item \textbf{Hard-coded unit prefactors.} Different magic factors appear across closely
    related classes: $\times 10$ in \textsc{ChargedTransport(Anisotropy)} and the
    \textsc{CurrentDensityAnisotropy} aux; $\times 100$ in \textsc{SideTotCrnt};
    $\times 10\,000$ in \textsc{ParamDiffusion}, \textsc{NeumannLi/En}, \textsc{MiecSE*},
    and the \textsc{CurrentDensity} aux. These are not interchangeable and downstream input
    files must be aware of which class they paired with.
  \item \textbf{Latent / suspected bugs spotted.}
    \begin{itemize}
      \item \textsc{ADRankTwoTensorNormalToBoundary} constructs \texttt{p1} via the comma
        operator (returns only the last component), so the auxiliary value is not the
        intended $\nvec^{\!\top}\bm{T}\nvec$.
      \item \textsc{ADCoupledButlerVolmerForce} returns the BV expression \emph{without}
        an explicit $\psi$ multiplier; verify against intent.
      \item \textsc{ChemicalEnergy} returns the bulk energy density $w(u-1)^2 u^2$ rather
        than its derivative; if used as a phase-field driving term this is
        mathematically inconsistent.
      \item \textsc{ADCurrentFluxEEL} default $\F=964523$ (vs.\ $96485$).
    \end{itemize}
  \item \textbf{Duplicated implementations.}
    \begin{itemize}
      \item \textsc{MiecSEAgC} and \textsc{MixAgC} are identical except for the registered
        class name.
      \item \textsc{LiPotInSE} and \textsc{LiPotMiec} are identical.
      \item \textsc{NernstPlanckConvection} (legacy) and \textsc{NernstPlanckConvectionScaled}
        differ only by an explicit \texttt{scale}.
    \end{itemize}
  \item \textbf{AD vs.\ legacy duplicates.} The legacy \textsc{Conduction}/\textsc{Migration}
    /\textsc{Kinetics}/\textsc{NernstPlanckConvection*} kernels predate their AD analogues
    \textsc{ADMigrationDiffusion}, \textsc{ADNernstPlanckConvection}, and the
    \textsc{ADCoupledButlerVolmerForce}; both families are still registered.
  \item \textbf{Non-standard registration.} \textsc{CZMComputeVariableJumpSmallStrain} is
    registered to \texttt{TensorMechanicsApp}, not \texttt{ecmApp}; consumers must include
    the upstream module.
  \item \textbf{AD sensitivity loss.} A handful of classes call \texttt{raw\_value} on AD
    properties (\textsc{ChargedTransport(Anisotropy)}, the flux aux kernels, \textsc{SideTotCrnt})
    or use non-AD \texttt{getMaterialProperty} (\textsc{ADThermalFlux}); the Jacobian therefore
    misses derivatives through those properties.
\end{enumerate}

\end{document}
